% Variant B only: included by main-smoothed.tex through the JRSmoothedVariant switch.
\section{Surface-Tension Pocket Smoothing}
\label{sec:pocket-smoothing}

\textit{Why the walls were chipped.} The exported solid is a height field over a fixed planform: faces $z=\pm t(x,y)/2$, with side walls on the outline and on every through-hole. Every weight-saving feature (isogrid pockets, lattice dimples, fullers, flutes, relief) is a removal depth $D=t_{\mathrm{base}}-t$. Before this release that depth reached the CAD through the $0.25\,\mathrm{mm}$ thickness raster by nearest-cell lookup, so most pocket walls were vertical steps along a staircase. A constrained Delaunay triangulation (CDT) of about $0.7\,\mathrm{mm}$ triangles cut across those steps, so the walls render as chipped (Figure~\ref{fig:pocket-bw}b, left). Mesh smoothing (\code{docs/paper/smoothing-ladder.pdf}) cannot remove this: it pins every vertex on an edge sharper than $30^\circ$, which is where the chips are.

\textit{Method.} We smooth the removal depth itself, before meshing, as a liquid would. Each pocket is a basin: a connected component of $D>0.02\,\mathrm{mm}$ (a watershed), merged across gaps narrower than $r/2$. Inside its basin the depth evolves for time $T=r^2/2$ by level-set curvature motion with a per-basin pressure,
\begin{equation}
\label{eq:pocket-flow}
D_t=|\nabla D|\,\nabla\!\cdot\!\frac{\nabla D}{|\nabla D|}+\lambda_b\,|\nabla D|,
\end{equation}
where $\lambda_b$ keeps the removed volume of basin $b$ constant (the volume of oil in the pit). Every depth contour then shortens like a two-dimensional droplet of fixed area under line tension: chips and spikes, which have the highest curvature, go first; straight rib edges stay straight and corners round to a radius of about $r$; and walls keep their steepness, so ribs keep their width. No basin may end deeper than its own deepest original point, so the volume constraint widens a thin groove instead of drilling it toward the minimum web. The planform, through-holes, grind faces, bore collar, raised features, brand pad and a $1\,\mathrm{mm}$ rim of unpocketed material along the outline and holes never move.

The flow runs on a $0.125\,\mathrm{mm}$ grid (about 3\,s for this design); the exporter reads $t=t_{\mathrm{base}}-D$ by bilinear interpolation and adds CDT vertices in three rows along every wall (depth contours at 8, 50 and 92\,\% of the local wall height) and in pairs straddling the grind boundary, so triangle edges follow each wall instead of cutting across it. Equation~(\ref{eq:pocket-flow}) is the $a=0$ case; $a=1$, a liquid-film flow, is compared below. The switch is \code{JARAMBIT_POCKET_SMOOTH} (\code{crates/crescent/src/pocket_smooth.rs}); \code{on} selects $r=0.75\,\mathrm{mm}$ and $a=0$. It is off by default, and with it unset every export is byte-identical to the previous release. An enabled setting is part of the certificate identity, so smoothed and as-exported evidence never share a key.

\begin{figure*}[t]
\centering\small
\begin{minipage}[c]{0.50\textwidth}\centering
\includegraphics[width=\linewidth]{pocket/hero.png}\\[1pt]{\footnotesize (a) Retained plate responses of the smoothed design}
\end{minipage}\hfill
\begin{minipage}[c]{0.47\textwidth}\centering
\includegraphics[width=0.49\linewidth]{pocket/bw-guard-raw.png}\hfill\includegraphics[width=0.49\linewidth]{pocket/bw-guard-ours.png}\\[1pt]
\includegraphics[width=0.49\linewidth]{pocket/bw-handle-raw.png}\hfill\includegraphics[width=0.49\linewidth]{pocket/bw-handle-ours.png}\\[1pt]
{\footnotesize (b) Raw export (left) and $r=0.75\,\mathrm{mm}$ (right): guard, handle}
\end{minipage}
\vspace{-1mm}
\caption{Surface-tension pocket smoothing of \code{w406t3_r01_p10_c16} at $r=0.75\,\mathrm{mm}$ ($F=80.370405$). (a)~The five retained plate responses of the smoothed design (bottom to top: unit lateral bending, chord torsion, ring traction, in-plane tip pry and the multi-case envelope), each normalized by its own area-weighted 99th percentile, with their normalized envelope on top; colour is a within-test rank, not a safety factor. (b)~Guard (isogrid pockets) and handle (flutes, lattice dimples), raw and smoothed, under the ladder's guard camera with $2.5\times$ thickness; the outline, holes and grinds are unchanged, and pockets keep their volume, depth and rib widths.}
\label{fig:pocket-bw}
\label{fig:pocket-fields}
\end{figure*}

\textit{Settings and fitness.} We exported the design at $r=0.5$, $0.75$, $1.0$, $1.5$ and $2.0\,\mathrm{mm}$ with $a=0$, and at $r=0.75\,\mathrm{mm}$ with $a=1$. Every export was re-assessed with the complete Stage-2 evaluator, so all twenty objectives, including the STEP solid, are measured on the smoothed geometry; mesh smoothing, by contrast, loses the 5-point CAD/B-rep objective (75.06--75.34 on the ladder). Table~\ref{tab:pocket-score} gives the results together with an unsmoothed re-run of the same evaluator, which reproduces the retained $80.299390$ exactly.

\begin{table}[t]
\centering\small
\caption{Complete Stage-2 re-assessment of \code{w406t3_r01_p10_c16} after pocket smoothing. $F$ is \code{W_ULTRALIGHT_v1} out of 100; the tip columns are the 3D P2 solid-tip safety factor at $1.82\,\mathrm{N\,m}$ (nodal maximum and average); lateral is the hunter capacity at certification and strict meshes.}
\label{tab:pocket-score}
\resizebox{\columnwidth}{!}{%
\begin{tabular}{@{}lrrrrrr@{}}
\toprule
Setting & $F$ & $\Delta F$ & Mass (g) & Tip SF & Lateral (N) & Battery \\
\midrule
As exported (re-run) & 80.299390 & -- & 83.2231 & 0.6823 / 0.7569 & 181.5 / 182.2 & 14/18 \\
$r=0.5$, $a=0$ & 80.164307 & $-0.135083$ & 83.2178 & 0.6353 / 0.6645 & 181.7 / 182.4 & 15/18 \\
$\boldsymbol{r=0.75}$, $a=0$ (chosen) & \textbf{80.370405} & $\mathbf{+0.071015}$ & 83.2089 & 0.7024 / 0.7511 & 181.8 / 182.5 & 14/18 \\
$r=1.0$, $a=0$ & 80.308595 & $+0.009205$ & 83.2070 & 0.6811 / 0.7237 & 181.9 / 182.5 & 14/18 \\
$r=1.5$, $a=0$ & 80.158645 & $-0.140745$ & 83.1869 & 0.6275 / 0.6751 & 182.2 / 182.7 & 14/18 \\
$r=2.0$, $a=0$ & 78.275029 & $-2.024361$ & 83.2001 & unmeasured & 182.3 / 182.8 & 13/18 \\
$r=0.75$, $a=1$ & 80.161219 & $-0.138171$ & 83.2629 & 0.6349 / 0.6752 & 182.2 / 182.7 & 14/18 \\
\bottomrule
\end{tabular}%
}
\end{table}

\textit{Results.} Smoothing leaves the design's engineering responses essentially where they were (Figure~\ref{fig:pocket-fields}a). The removed volume of every basin is preserved to rounding, and the exported volume changes by $-3.19\,\mathrm{mm^3}$ ($-0.017\,\%$) at $r=0.75\,\mathrm{mm}$. Mass stays at $83.19$--$83.26\,\mathrm{g}$ against $83.2231\,\mathrm{g}$ as exported; the lateral capacity rises by $0.2$--$0.8\,\mathrm{N}$; the ring, body-twist and fatigue margins move by less than $3\,\%$; the failure-sequence ratio rises from $1.29$ to $1.45$--$1.62$; and the CAD/B-rep objective keeps its five points in every case.

What moves $F$ is the 3D solid-tip submodel. Its nodal-maximum safety factor changes non-monotonically with $r$: $0.6353$, $0.7024$, $0.6811$ and $0.6275$ for $r=0.5$ to $1.5\,\mathrm{mm}$, against $0.6823$ as exported. At $r=2.0\,\mathrm{mm}$ the tip solve produced no observation, which the evaluator scores as zero. We read these changes as the sensitivity of a nodal maximum near the fuller end to the local triangulation, not as a structural trend (the two solvers agree within $0.08\,\%$). The setting $r=0.75\,\mathrm{mm}$, chosen visually, is also the best measured: $F=80.370405$, $0.071015$ above the as-exported design. It is the design shown on the first page (Figure~\ref{fig:winner-cover}) and in Figure~\ref{fig:cover}.

Larger $r$ makes pockets rounder and more alike. At $r=2.0\,\mathrm{mm}$ small pockets merge into rounded blobs, and the film variant ($a=1$) spreads depth sideways into the ribs; neither improves $F$. The chosen export has about 40\,\% more triangles than the raw one (STEP $86.7$ against $53.4\,\mathrm{MB}$), and its Stage-2 assessment took about three hours instead of about 35 minutes, mostly in STEP verification and the larger tip solve ($177{,}984$ against $91{,}008$ tetrahedra).

\textit{Scope.} The smoothing changes interior pocket geometry only and adds no qualification or physical observation. The Individual~2000 plate fields of Section~\ref{sec:fitness} are retained evidence of a different design and are not smoothed. Receipts, scorecards and every exploratory setting are in \code{experiments/pocket-surface-tension-smoothing/}.
